\section[Una sucesión exacta fundamental]{Una sucesión exacta fundamental. Descenso por morfismos con
fibras relativamente conexas}
\label{IX.6}
\marginpar{253}

\begin{theoreme}
\label{IX.6.1}
Sean~$S$ el espectro de un anillo artiniano~$A$ con cuerpo residual~$k$,
~$\overline{k}$ una clausura algebraica de~$k$,~$X$ un~$S$-preesquema,
~$X_0=X\otimes_A k$,~$\overline{X}_0=X\otimes_A\overline{k}$,
~$\overline{a}$ un punto geométrico
\ifthenelse{\boolean{orig}}{de~$\overline{X}$}{de~$\overline{X}_0$},
~$a$ su imagen en~$X$ y~$b$ su imagen en~$S$. Supongamos que~$X_0$ es
cuasicompacto y \emph{geométricamente conexo} sobre~$k$ (N.B.: si~$X$ es
propio sobre~$S$, esto significa que~$\H^0(X_0,\cal{O}_{X_0})$ es un anillo
artiniano \emph{local} cuyo cuerpo residual es \emph{radicial} sobre~$k$).
Entonces la sucesión de homomorfismos canónicos
$$
e\mto \pi_1(\overline{X}_0,\overline{a})\to
\pi_1(X,a)\to\pi_1(S,b)\to e
$$
es exacta, y se tiene
$$
\pi_1(S,b)\isomfrom \pi_1(k,\overline{k})=\text{grupo de Galois
de~$\overline{k}$ sobre~$k$.}
$$
\end{theoreme}

Como los grupos fundamentales no cambian al eliminar los elementos
nilpotentes, podemos suponer que~$A=k$, lo cual hace ya evidente el último
isomorfismo. Sea~$k'$ la clausura separable de~$k$ en~$\overline{k}$,
consideremos~$X'=X\otimes_k k'$ y sea~$a'$ la imagen de~$\overline{a}$
en~$X'$. Se tiene una sucesión de homomorfismos canónicos
$$
e\mto \pi_1(\overline{X}_0,\overline{a})\to
\pi_1(X',a')\to\pi_1(S',b')\to e
$$
(donde~$S'=\Spec(k')$). Por último, el diagrama
$$
\xymatrix{ S & X \ar[l] & \overline{X}_0 \ar[l] \\ S' \ar[u] &
X'\ar[u] \ar[l] & \overline{X}_0 \ar[l] \ar@{=}[u]. }
$$
da un homomorfismo canónico de esta sucesión en la correspondiente a~$X/k$.
Por otra parte, se ve que este homomorfismo de sucesiones de grupos es un
isomorfismo, como resulta de~\Ref{IX.4.11}. Quedamos, pues, reducidos a probar
que la segunda sucesión es exacta, es decir, podemos suponer que~$k$ es
\emph{perfecto}. Sean entonces~$k_i$ las subextensiones galoisianas finitas
de~$k$ en~$\overline{k}$, pongamos~$X_i=X\otimes_k k_i$ y sea~$a_i$ la
imagen de~$\overline{a}$ en~$X_i$. Dejamos al lector
\ifthenelse{\boolean{orig}}
{\ignorespaces}
{la tarea}
de comprobar que el homomorfismo natural
$$
\pi_1(\overline{X}_0,\ifthenelse{\boolean{orig}}{a}{\overline{a}})
\isomto \varprojlim\pi_1(X_i,a_i)
$$
\marginpar{254}%
es un isomorfismo. Esto significa simplemente que todo recubrimiento étale
de~\ifthenelse{\boolean{orig}}{$\overline{X}$}{$\overline{X}_0$} procede de
un recubrimiento étale de algún~$X_i$, y que este último es esencialmente
único, salvo pasar a un~$X_j$,~$j\geq i$.
Por otra parte, sea~$\pi_i$ el grupo de Galois de~$k_i$ sobre~$k$, es decir,
el grupo opuesto al grupo de los~$S$-automorfismos de~$S_i=\Spec(k_i)$.
Como el funtor~$S'\mto X\times_S S'$ de los recubrimientos étale de~$S$ en
los recubrimientos étale de~$X$ es plenamente fiel~\eqref{IX.3.4}, se
sigue que~$\pi_i$ es también isomorfo al grupo opuesto al grupo de los
~$X$-automorfismos del recubrimiento principal conexo~$X_i$ de~$X$. Resulta,
por tanto, de~\Ref{V.6.13} que se tiene una sucesión exacta
$$
e\mto \pi_1(X_i,a_i)\to \pi_1(X,a)\to\pi_i\to e.
$$
Pasando al límite proyectivo sobre~$i$ en estas sucesiones exactas, se obtiene
una sucesión exacta (pues estamos en la categoría de los grupos de tipo
galoisiano), que no es otra que la sucesión considerada en~\Ref{IX.6.1}.
Esto termina la demostración.

La traducción de la exactitud por la derecha en~\Ref{IX.6.1} al lenguaje
geométrico es la siguiente:

\begin{corollaire}
\label{IX.6.2}
Con las notaciones precedentes, sea~$X'$ un recubrimiento étale de~$X$ y sea
~$\overline{X'}_0$ el recubrimiento étale correspondiente de
~$\overline{X}_0$. Las condiciones siguientes son equivalentes:
\begin{enumerate}
\item[(i)] Existe un~$S'$ étale sobre~$S$ y un~$X$-isomorfismo
$X'\isomto X\times_S S'$ (entonces~$S'$ queda determinado salvo isomorfismo
único en virtud de~\Ref{IX.3.4}).
\item[(ii)] $\overline{X}_0'$ está completamente descompuesto sobre
~$\overline{X}_0$.

Si~$X'$ es conexo, estas condiciones equivalen también a:
\item[(ii~bis)] $\overline{X}_0'$ posee una sección sobre~$\overline{X}_0$.
\end{enumerate}
\end{corollaire}

(N.B.: Este último complemento es esencial; la equivalencia de~(i) y~(ii)
solo significa que~$\pi_1(S,b)$ es el grupo cociente de~$\pi_1(X,a)$ por el
subgrupo
\marginpar{255}
normal cerrado engendrado por la imagen de
~$\pi_1(\overline{X}_0,\overline{a})$, y no por esta imagen misma). Bajo las
condiciones precedentes, diremos que~$X'$ es un recubrimiento
\emph{geométricamente trivial}
\index{geométricamente trivial (recubrimiento)|hyperpage}%
de~$X$.

\begin{remarque}
\label{IX.6.3}
En el enunciado~\Ref{IX.6.1} no puede sustituirse~$\overline{k}$ por una
extensión algebraicamente cerrada cualquiera de~$k$, ni siquiera si se supone
ya que~$k$ es algebraicamente cerrado. En otras palabras, no es cierto en
general que, si~$X$ es un esquema algebraico conexo sobre un cuerpo
algebraicamente cerrado~$k$, su grupo fundamental permanezca inalterado al
sustituir~$k$ por una extensión algebraicamente cerrada. Esto es falso ya, por
ejemplo, en característica~$p>0$ para la recta afín sobre~$k$, a causa de los
fenómenos de «ramificación superior» en el punto del infinito, que introducen
una estructura «continua» en el grupo fundamental. Veremos, sin embargo, en el
exposé siguiente que tales fenómenos no pueden producirse si~$X$ es
\emph{propio} sobre~$k$. Demostraremos también por métodos trascendentes que
ocurre lo mismo si~$k$ es de característica cero.
\end{remarque}

\begin{corollaire}
\label{IX.6.4}
Supongamos que~$a$ está localizado en un punto~$x\in X$ racional sobre~$k$ (o,
más generalmente, cuyo cuerpo residual es radicial sobre~$k$). Entonces la
sucesión exacta~\Ref{IX.6.1} es escindida.
\end{corollaire}

Podemos suponer~$S=\Spec(k)$. Si~$x$ es racional sobre~$k$, corresponde a una
sección~$S\to X$ de~$X$ sobre~$S$, que lleva~$b$ en~$a$ y define un
homomorfismo~$\pi_1(S,b)\to\pi_1(X,a)$ que es la escisión buscada. Si
~$\kres(x)$ es radicial sobre~$k$, nos reducimos al caso precedente mediante
la extensión de base~$\Spec(\kres(x))\to\Spec(k)$.

\begin{theoreme}
\label{IX.6.5}
Sean~$f\colon X\to S$ un morfismo propio y sobreyectivo de presentación
finita, con fibras geométricamente conexas,~$X'$ un preesquema de presentación
finita y propio sobre~$X$,~$s$ un punto de~$S$,~$F=X_s$ la fibra de~$X$ en
~$s$, y~$F_1'$ una componente conexa de la fibra~$F'=X'_s$ de~$X'$ en~$s$.
Para que existan un entorno abierto~$X_1'$ de~$F_1'$ en~$X'$, un
~$S$-esquema étale~$S_1'$ y un~$X$-isomorfismo
$X_1'\isomto S_1'\times_S X$, es necesario y suficiente que~$X'$ sea étale
sobre~$X$ en los puntos de~$F_1'$ y que~$F_1'$ sea un recubrimiento
geométricamente trivial de~$F$.
\end{theoreme}

La
\marginpar{256}
necesidad de la condición es trivial, por lo que resta probar la suficiencia.
Nos reducimos fácilmente al caso en que~$S$ es noetheriano. Consideremos la
factorización de Stein~$X\to T\to S$ de~$f$, donde~$T$ es el espectro del
álgebra~$f_*(\cal{O}_X)$ sobre~$S$. Como las fibras de~$X$ sobre~$S$ son
geométricamente conexas y~$f$ es sobreyectivo, el morfismo~$T\to S$ es finito,
sobreyectivo y radicial; por consiguiente, según~\eqref{IX.4.10}, todo~$T'$
étale sobre~$T$ procede por imagen inversa de un~$S'$ étale sobre~$S$.
Esto nos reduce a
\ifthenelse{\boolean{orig}}
{probar~\Ref{IX.6.5}.}
{probar~\Ref{IX.6.5}}
reemplazando allí~$S$ por~$T$, es decir, al caso en que suponemos
~$f_*(\cal{O}_X)=\cal{O}_S$. Consideremos entonces la factorización de Stein
~$X'\to S'\to S$ del morfismo propio~$h\colon X'\to S$, donde~$S'$ es el
espectro del álgebra~$h_*(\cal{O}_{X'})$. Los morfismos~$X'\to X$ y
~$X'\to S'$ definen un morfismo canónico
$$
X'\to X\times_S S',
$$
y nuestra afirmación está contenida en la siguiente:
